\section{Additional results for \texorpdfstring{Section~\ref{sec:env}}{Section 14}: environments and selective sensing}\label{sup:env}

This section collects statements, examples and remarks of \cref{sec:env} that the main text cites without stating in full. Proofs are in \cref{app:env}.

\subsection{Summaries of the past}\label{sup:env-summaries-of-the-past}

\begin{lemma}[Summaries of the past do no better than sensors of the present]\label{lem:env-markov}
Let $\eta$ be stationary Gauss--Markov and $\mathsf Y=\mathsf S_{\mathrm m}\eta_{t+h}$, with a selector $\mathsf S_{\mathrm m}$ of full row rank, and let $\mathcal S$ be a summary of the past that is jointly Gaussian with the process and conditionally independent of $(\eta_{t+u})_{u>0}$ given $(\eta_{t'})_{t'\le t}$, for example a linear function of the past plus independent noise, with $\dim\mathcal S\le\kappa$ and $I(\mathcal S;(\eta_{t'})_{t'\le t})\le I_{\max}$. Then $I(\mathcal S;\mathsf Y)$ is at most the value of \cref{thm:env-canonical}(b) for $\mathsf X=\eta_t$.
\end{lemma}
% src: H5-L3 (selector renamed S_m, critic 2.7)

Proof in \cref{app:env}; checked in \cref{chk:env-markov}. For a fully observed Markov exterior the past-future canonical variates are therefore those of the pair $(\eta_t,\eta_{t+h})$, with $\Sigma_{\mathsf X}=P_0$ and $\Sigma_{\mathsf X\mathsf Y}=P_0(F^h)^\top\mathsf S_{\mathrm m}^\top$. An agent whose raw input is an observation process applies \cref{thm:env-canonical} with $\mathsf X$ a window of past observations and $\mathsf Y$ the observation at the horizon, the linear form of the past-future bottleneck of \citet{creutzig2009past}. Predictive information, the mutual information between past and future of a time series, is defined by \citet{bialek2001predictability}, and canonical variables of past and future give Markovian representations \citep{akaike1975markovian}, a work cited from its record.

\subsection{Diffusive exteriors}\label{sup:env-diffusive-exteriors}

\begin{proposition}[Diffusive exteriors: long horizons read coarse variables]\label{prop:env-diffusive}
Let the exterior be a field on a connected graph with at least two vertices and symmetric weighted Laplacian $L_{\mathrm w}$, $F=(1-\varepsilon_{\mathrm w})I-\mathsf d_{\mathrm w}L_{\mathrm w}$ with $\varepsilon_{\mathrm w}\in(0,1)$, $\mathsf d_{\mathrm w}>0$ and $1-\varepsilon_{\mathrm w}-\mathsf d_{\mathrm w}\lambda_{\max}(L_{\mathrm w})>0$, and $Q=Q_0I$. Then the exterior is reversible, and at every horizon its canonical variates are, up to normalization, the projections on an orthonormal basis of Laplacian eigenvectors, ordered by increasing Laplacian eigenvalue: the slowest is the global average, with $\tau_1=-1/\ln(1-\varepsilon_{\mathrm w})$, and the next is the projection on an eigenvector of the algebraic connectivity $\lambda_2(L_{\mathrm w})$, the second-smallest Laplacian eigenvalue, with $\tau_2=-1/\ln(1-\varepsilon_{\mathrm w}-\mathsf d_{\mathrm w}\lambda_2(L_{\mathrm w}))$. Under a budget $I_{\max}>0$, the budget-optimal sensor of \cref{thm:env-canonical}(b) with $\kappa$ channels, for $\mathsf X=\eta_t$ and $\mathsf Y=\eta_{t+h}$, reads the lowest-frequency modes, at most $\kappa$ of them and no more as $h$ grows (modes of equal Laplacian eigenvalue share one budget threshold), and only the global average once $h>I_{\max}/(\tau_2^{-1}-\tau_1^{-1})$.
\end{proposition}
% src: H5-P2

Proof in \cref{app:env}. In a diffusive exterior slow and coarse coincide: a long-horizon agent senses spatial averages and long-wavelength modes, which is the precise sense in which a slower agent has a more abstract environment. Whether the modes that a Gaussian coarse-graining retains are those a horizon-$h$ agent senses is a question for \cref{sec:ren}; in reversible diffusive exteriors both are the slow Laplacian modes.

\subsection{Expected free energy and preferences}\label{sup:env-expected-free-energy-and-preferences}

\begin{remark}[The attention expected free energy of the framework is horizon-free]\label{rem:env-efe-horizon-free}
With world channels only and no preference over them, \cref{prop:fw-efe-communication}(a) scores attention by $-I(\eta;y_{\mathcal J})$, information about the present state. For a linear sensor of $\eta_t$ this is $\frac12\sum_{\mathsf j}\ln(1-\mathsf g_{\mathsf j})$ by \cref{lem:env-information}(i), a function of the captured fractions alone: under a binding Shannon budget every sensor that spends the budget has the same epistemic value, whatever it reads. Under a gain budget $\lvert f\rvert\le1$ on one channel $s=f^\top\eta_t+v$, the value $\frac12\ln(1+f^\top P_0f)$ is largest along the top eigenvector of $P_0$, the direction of largest variance. For two independent modes with $\lambda(F)=(0.95,0.5)$ and innovation variances $(0.01,1)$ the stationary variances are $(0.103,1.333)$, so this sensor reads the fast mode, while the canonical sensor reads the slow mode at every horizon (\cref{chk:env-control}). Selection by predictive information needs a horizon in the epistemic term, as in \cref{def:env-attention-objective}, or a pragmatic term.
\end{remark}
% src: H5-R2

\begin{remark}[The literal two-step expected free energy rewards complementary sensors]\label{rem:env-two-step}
Score the two-step policy (sense at $t$, act, observe $o$ at $t+h$) by the expected free energy of \cref{prop:fw-efe-communication}, minus epistemic minus pragmatic value, with the path $\eta_{t:t+h}$ as hidden state, $(s,o)$ as outcomes and $R^{\mathrm o}\succ0$. Since $s$ depends on the path through $\eta_t$ and, given $s$, the act is known and $o$ depends on the path through $\eta_{t+h}$, the epistemic term is $I(\eta_t;s)+I(\eta_{t+h};o\mid s)$, and with $\tilde o=H^{\mathrm o}\eta_{t+h}+v^{\mathrm o}$ Gaussian conditioning gives
\[
I(\eta_{t+h};o\mid s)=\tfrac12\ln\det\bigl(H^{\mathrm o}\Cov(\eta_{t+h}\mid s)(H^{\mathrm o})^\top+R^{\mathrm o}\bigr)-\tfrac12\ln\det R^{\mathrm o},
\]
a constant minus $I(s;\tilde o)$. Under a binding budget the first term is fixed, and the epistemic term prefers the sensor least informative about the future outcome: a reading is worth less when the outcome will reveal the same thing later. This opposes the pragmatic term of \cref{thm:env-control}(a), which dominates when preferences are precise. The predictive-information criterion of \cref{thm:env-canonical} follows from the expected free energy only when the hidden state of the agent's model for this decision is the state at the horizon and the later outcome $o$ is not scored in the epistemic term: with hidden state $\eta_{t+h}$ and outcomes $(s,o)$ the epistemic term is $I(s;\eta_{t+h})+I(\eta_{t+h};o\mid s)$, which by the display equals $I(s;\eta_{t+h}\mid\tilde o)$ plus a constant and, as a function of the sensor, differs from the predictive information unless $H^{\mathrm o}F^h=0$. We therefore declare the predictive information as the objective of \cref{def:env-attention-objective} and claim no more.
\end{remark}
% src: H5-R3

\begin{example}[A preference on a fast variable reverses the selection]\label{ex:env-fast-preference}
Take two independent modes, slow ($\lambda(F)=0.95$, innovation variance $0.1$) and fast ($\lambda(F)=0.5$, innovation variance $1$, stationary variance $\mathsf p_{\mathrm f}=4/3$), $\mathsf S_{\mathrm m}=I$, one channel with $I_{\max}=1$ nat, and $\Upsilon_{\mathrm{ctl}}=\diag(0,w_{\mathrm p})$ with $w_{\mathrm p}\in\{4,40\}$, so that the preferred and controllable outcome is the fast mode. With $x$ the squared sine of the sensor's angle from the slow mode in whitened coordinates, $\mathsf g=1-e^{-2}$, $\sigma_{\mathrm s}^2=0.95^{2h}$ and $\sigma_{\mathrm f}^2=0.25^h$, up to a constant
\[
G^{\mathrm{att}}_h=\tfrac12\ln\bigl(1-\mathsf g\sigma_{\mathrm s}^2+\mathsf gx(\sigma_{\mathrm s}^2-\sigma_{\mathrm f}^2)\bigr)-\tfrac12\mathsf gw_{\mathrm p}\mathsf p_{\mathrm f}\sigma_{\mathrm f}^2x,
\]
which is concave in $x$, so the optimum reads a pure mode. With $w_{\mathrm p}=4$ the slow mode is read at every $h\in\{1,\dots,8\}$, by its epistemic value alone; with $w_{\mathrm p}=40$ the fast mode is read at $h=1,2$ and the slow mode from $h=3$ on (\cref{chk:env-control}). The pragmatic optimum reads the fast mode and the epistemic optimum the slow mode at every $h$. At long horizons this agent reads a variable its preferences ignore, because $\mathsf S_{\mathrm m}=I$ makes it curious about everything; with the model target restricted to the fast mode the slow mode has no value (\cref{prop:env-selection}(ii)).
\end{example}
% src: H5-X1; horizons 1 to 8 rechecked by the editor

\subsection{Selection principle and the cell model}\label{sup:env-selection-principle-and-the-cell-model}

\begin{proposition}[Selection principle]\label{prop:env-selection}
Let the exterior split into contents $\eta=(\eta^{(c)})_{c\in\mathcal M}$ as in \cref{def:abs-multiplex}, let the agent hold channels $s^{\mathrm{old}}$ jointly Gaussian with $\eta$, and add a channel $s^{(c)}=H^{(c)}\eta^{(c)}_t+v^{(c)}$ with independent nondegenerate Gaussian noise. Under $G^{\mathrm{att}}_h$ its value is the sum of the epistemic increment $\mathsf V_{\mathrm{ep}}=I(s^{(c)};\mathsf S_{\mathrm m}\eta_{t+h}\mid s^{\mathrm{old}})$ and the pragmatic increment $\mathsf V_{\mathrm{pr}}=\frac12\tr\bigl(\Upsilon_{\mathrm{ctl}}[\Cov(\eta_{t+h}\mid s^{\mathrm{old}})-\Cov(\eta_{t+h}\mid s^{\mathrm{old}},s^{(c)})]\bigr)$.
\begin{enumerate}
\item[(i)] (Coupling and predictability.) If $\Cov(s^{(c)},\eta_{t+h}\mid s^{\mathrm{old}})=0$, both increments vanish. This covers no coupling, $H^{(c)}=0$, and a content uncorrelated at lag $h$ with the exterior given what the agent already reads. Both increments are continuous in $H^{(c)}$.
\item[(ii)] (Relevance.) $\mathsf V_{\mathrm{pr}}=0$ if and only if $\Cov(s^{(c)},\mathbb P_{\mathrm o}H^{\mathrm o}\eta_{t+h}\mid s^{\mathrm{old}})=0$, which holds when the content neither enters the preferred outcomes nor correlates with what does, or when it enters them only along directions that acts cannot compensate and, given $s^{\mathrm{old}}$, is uncorrelated with what enters the compensable directions; $\mathsf V_{\mathrm{ep}}=0$ if and only if $\Cov(s^{(c)},\mathsf S_{\mathrm m}\eta_{t+h}\mid s^{\mathrm{old}})=0$.
\item[(iii)] (Too fast.) If content $c$ is a stationary Gauss--Markov block with transition $F^{(c)}$ and stationary covariance $P^{(c)}_0$, independent of the other contents and of the inputs of $s^{\mathrm{old}}$, then for every gain
\[
\mathsf V_{\mathrm{ep}}\le-\tfrac12\sum_{\mathsf j}\ln\bigl(1-\sigma^{(c)}_{\mathsf j}(h)^2\bigr),\qquad\mathsf V_{\mathrm{pr}}\le\tfrac12\Vert\Upsilon_{\mathrm{ctl}}\Vert\tr\bigl(F^{(c)h}P^{(c)}_0(F^{(c)h})^\top\bigr),
\]
where $\sigma^{(c)}_{\mathsf j}(h)$ are the canonical correlations of the block at lag $h$. They are at most $\mathsf c^{(c)}\Vert F^{(c)h}\Vert$, with $\mathsf c^{(c)}=\Vert(P^{(c)}_0)^{1/2}\Vert\Vert(P^{(c)}_0)^{-1/2}\Vert$, so both bounds vanish geometrically in $h$.
\item[(iv)] (Constant at the agent's scale.) A content whose innovation is small against the sensor noise yields vanishing information per reading in steady state, whatever its persistence (\cref{lem:env-constant}).
\end{enumerate}
\end{proposition}
% src: H5-P3

Proof in \cref{app:env}. Prediction fails for variables that relax much faster than the horizon and, per reading, for variables that barely change over the agent's lifetime. Curiosity values every predictable variable in the model, so selection by relevance needs a model target restricted to relevant variables or a budget under which pragmatic value dominates.

\begin{lemma}[Variables constant at the agent's scale]\label{lem:env-constant}
Let a scalar exterior variable follow $\eta^{\mathrm g}_{t+1}=F_{\mathrm g}\eta^{\mathrm g}_t+w^{\mathrm g}_t$ with $\lvert F_{\mathrm g}\rvert\le1$ and innovation variance $Q_{\mathrm g}>0$, and be read as $s_t=\eta^{\mathrm g}_t+v_t$ with $v_t\sim\Ncal(0,R^s)$, $R^s>0$. From any prior, the predicted variance $P^-_t$ converges to the positive root $P^-_\infty$ of $x^2+x(R^s(1-F_{\mathrm g}^2)-Q_{\mathrm g})-Q_{\mathrm g}R^s=0$, and the information that a new reading carries given the past readings, $\frac12\ln(1+P^-_t/R^s)$, converges to
\[
\tfrac12\ln\bigl(1+P^-_\infty/R^s\bigr)\le\tfrac12\ln\Bigl(1+\frac{Q_{\mathrm g}}{R^s}+\Bigl(\frac{Q_{\mathrm g}}{R^s}\Bigr)^{1/2}\Bigr);
\]
the information about $\eta^{\mathrm g}_{t+h}$ is no larger. For a constant ($Q_{\mathrm g}=0$, $F_{\mathrm g}=1$) with prior variance $P_{\mathrm{pr}}$, $n_{\mathrm r}$ readings carry $\frac12\ln(1+n_{\mathrm r}P_{\mathrm{pr}}/R^s)$ nats in total, with a rate per reading that tends to zero.
\end{lemma}
% src: H5-L4

Proof in \cref{app:env}; checked in \cref{chk:env-constant}. A cell senses its ligand through a diffusive medium, whose molecular noise sets the coupling of the chemical channel.

\begin{proposition}[Volume-time average in an ideal ligand gas]\label{prop:env-concentration}
Let ligand molecules be independent Brownian particles in $\R^3$ with diffusion constant $D_{\mathrm m}$, at equilibrium with mean density $\bar c_{\mathrm L}$, so that their positions at each time form a Poisson process, and let a transparent ball $\mathbb B$ of radius $r_{\mathrm{cell}}$ and volume $\lvert\mathbb B\rvert$ average the number $\mathsf N_{\mathbb B}(t)$ of molecules inside it over $[0,T_{\mathrm{int}}]$: $\hat c=(\lvert\mathbb B\rvert T_{\mathrm{int}})^{-1}\int_0^{T_{\mathrm{int}}}\mathsf N_{\mathbb B}(t)\mathrm dt$. Then $\E\hat c=\bar c_{\mathrm L}$ and
\[
\frac{\Var\hat c}{\bar c_{\mathrm L}^2}=\frac{3}{5\pi D_{\mathrm m}r_{\mathrm{cell}}\bar c_{\mathrm L}T_{\mathrm{int}}}\Bigl[1-\frac53\Bigl(\frac{r_{\mathrm{cell}}^2}{\pi D_{\mathrm m}T_{\mathrm{int}}}\Bigr)^{1/2}+\mathfrak q_T\Bigr],
\]
with
\[
0\le\mathfrak q_T\le\frac{5}{3\sqrt\pi}\frac{r_{\mathrm{cell}}^2}{D_{\mathrm m}T_{\mathrm{int}}}+\frac{5}{18\sqrt\pi}\Bigl(\frac{r_{\mathrm{cell}}^2}{D_{\mathrm m}T_{\mathrm{int}}}\Bigr)^{3/2}.
\]
In an ideal solution a small change of chemical potential, in units of $k_{\mathrm B}T_{\mathrm K}$ with $T_{\mathrm K}$ the absolute temperature, equals the relative change of concentration, so the same number is, to first order in the relative error, the variance of the cell's estimate of its ligand's chemical potential in units of $(k_{\mathrm B}T_{\mathrm K})^2$.
\end{proposition}
% src: H5-P4 (leading term is the Berg--Purcell perfect-instrument formula, attributed through Mora--Wingreen Eq. 10; finite-time terms computed here)

Proof in \cref{app:env}; checked in \cref{chk:env-concentration}. The leading term is the known formula for a sphere that counts the molecules inside its volume without hindering their diffusion: \citet{mora2010limits} give $3/(5\pi D_{\mathrm m}r_{\mathrm{cell}}\bar c_{\mathrm L}T_{\mathrm{int}})$, in our symbols, as their Eq.~10 and attribute it to \citet{berg1977physics}, whose abstract states a least fractional error of about $(T_{\mathrm{int}}\bar c_{\mathrm L}r_{\mathrm{cell}}D_{\mathrm m})^{-1/2}$. We take the attribution of the constant from the former and know the latter from its abstract. The appendix rederives the leading term; the finite-time correction $-\frac53(r_{\mathrm{cell}}^2/(\pi D_{\mathrm m}T_{\mathrm{int}}))^{1/2}$ and the bound on $\mathfrak q_T$ are computed here. \Citet{bialek2005physical} bound the relative root-mean-square error of a single receptor of linear size $a_{\mathrm r}$ below by $(\pi D_{\mathrm m}a_{\mathrm r}\bar c_{\mathrm L}T_{\mathrm{int}})^{-1/2}$ (their Eq.~35 in the preprint, in our symbols); sensor models change the constant, not the scaling.

\begin{example}[A cell that senses chemical potential and not gravity]\label{ex:env-cell}
A cell of radius $r_{\mathrm{cell}}=5$~$\mu$m has an exterior with three components: the ligand concentration $\bar c_{\mathrm L}+\eta^{\mathrm L}_t+c'_{\mathrm L}(x,t)$, where $\eta^{\mathrm L}$ is a slow signal (an Ornstein--Uhlenbeck process with correlation time $\tau_{\mathrm L}=10$~s and a standard deviation of a few percent of $\bar c_{\mathrm L}$) and $c'_{\mathrm L}$ is the thermal fluctuation of \cref{prop:env-concentration}; the component $\eta^{\mathrm g}$ of the gravitational acceleration along the cell's axis, which changes with the organism's posture; and constant offsets. The cell's acts affect outcomes that its preferences weight through $\eta^{\mathrm L}_{t+h}$ only, at a horizon of seconds. Chemoreception is the volume-time average of \cref{prop:env-concentration} with $T_{\mathrm{int}}=1$~s, $D_{\mathrm m}=500$~$\mu$m$^2$/s and $\bar c_{\mathrm L}=1$~nM; the candidate gravity sensor is the equilibrium height of buoyant particles in a compartment of height $Z_{\mathrm{box}}=10$~$\mu$m. The computed values (\cref{chk:env-cell}) are these.
\begin{enumerate}
\item[(1)] Coupling of the chemical channel. The relative error is $1.13\times10^{-2}$ in the leading term and $1.01\times10^{-2}$ exactly: since $r_{\mathrm{cell}}^2/D_{\mathrm m}=0.05$~s, the finite-time factor on the variance is $0.81$ ($0.79$ without the remainder $\mathfrak q_T=0.020$). The values of this item use the leading term, which overstates the noise. A signal with standard deviation 5 (1) percent of $\bar c_{\mathrm L}$ has signal-to-noise ratio $19.7$ ($0.79$) per reading, captured fraction $0.952$ ($0.441$), and predictive information $0.755$ ($0.224$) nats at $h=1$~s and $0.069$ ($0.031$) nats at $h=10$~s; with the exact variance the ratios are $24.3$ ($0.97$) and the captured fractions $0.960$ ($0.493$). These ratios, fractions and information values treat the one-second average as a sample of $\eta^{\mathrm L}$ at the end of the window plus the thermal noise, with $h$ measured from the end of the window. The reading averages the signal as well: with $a=T_{\mathrm{int}}/\tau_{\mathrm L}=0.1$, the average of the Ornstein--Uhlenbeck signal has variance $2a^{-2}(a-1+e^{-a})=0.967$ times that of the signal and covariance $a^{-1}(1-e^{-a})=0.952$ times that variance with $\eta^{\mathrm L}$ at the end of the window, and the predictive information of the averaged reading is $0.651$ ($0.201$) nats at $h=1$~s and $0.064$ ($0.028$) nats at $h=10$~s in the leading term, and $0.664$ ($0.232$) and $0.065$ ($0.032$) nats with the exact variance.
\item[(2)] The thermal fluctuation is too fast. The autocorrelation of its volume average decays like $\lvert\mathbb B\rvert(4\pi D_{\mathrm m}u)^{-3/2}$ for $D_{\mathrm m}u\gg r_{\mathrm{cell}}^2$, and in the Gaussian approximation, with the exact variance of the reading, a one-second reading carries $2.0\times10^{-3}$ nats about the next one and $5.3\times10^{-7}$ nats about the one ten seconds ahead. These are computed values; the large-time forms are leading-order asymptotics without a remainder bound.
\item[(3)] Constant offsets ($\bar c_{\mathrm L}$, the magnitude of the acceleration) are learned once (\cref{lem:env-constant}).
\item[(4)] Coupling of gravity. A particle of buoyant mass $\mathsf M_{\mathrm b}$ has the truncated Boltzmann height law $\propto e^{-\mathsf M_{\mathrm b}\eta^{\mathrm g}\mathsf z/(k_{\mathrm B}T_{\mathrm K})}$ on $[0,Z_{\mathrm{box}}]$, with Fisher information $(\mathsf M_{\mathrm b}/(k_{\mathrm B}T_{\mathrm K}))^2\Var(\mathsf z)$ about $\eta^{\mathrm g}$, so that, with $g$ the gravitational acceleration and $x_{\mathrm g}=\mathsf M_{\mathrm b}gZ_{\mathrm{box}}/(k_{\mathrm B}T_{\mathrm K})$, one independent reading at $\eta^{\mathrm g}=g$ has signal-to-noise ratio $g^2$ times this Fisher information, which equals $1-(x_{\mathrm g}/2)^2/\sinh^2(x_{\mathrm g}/2)$, for changes of $\eta^{\mathrm g}$ of order $g$. This ratio is $x_{\mathrm g}^2/12$ to leading order for $x_{\mathrm g}\ll1$, increases with $x_{\mathrm g}$ and stays below one. For a protein of 50~kDa with buoyancy factor $0.27$ (partial specific volume $0.73$~cm$^3$/g) at 310~K, $x_{\mathrm g}=5.1\times10^{-7}$ and the ratio is $2.2\times10^{-14}$ per reading: about $4.5\times10^{13}$ independent readings give a ratio of one. A particle of radius 1~$\mu$m and density excess 500~kg/m$^3$ has $x_{\mathrm g}=48$, and sedimentation dominates thermal motion: its ratio is one to within $4\times10^{-18}$, the value for a height that is exponentially distributed with the scale height $Z_{\mathrm{box}}/x_{\mathrm g}$, so that in this model no single particle reaches the ratio $19.7$ of the 5 percent ligand signal per reading. The ratio measures small changes about $\eta^{\mathrm g}=g$. A reversal of $\eta^{\mathrm g}$ reflects the height law about the middle of the compartment, and the rule that reads the sign of $\eta^{\mathrm g}$ from the half in which the particle lies errs with probability $1/(1+e^{x_{\mathrm g}/2})$, which is $3.8\times10^{-11}$ at $x_{\mathrm g}=48$ and within $10^{-7}$ of one half for the protein.
\item[(5)] Relevance of gravity. In the model $\eta^{\mathrm g}$ is independent of $\eta^{\mathrm L}$ and absent from the preferred outcomes, so its pragmatic increment is zero (\cref{prop:env-selection}(ii)), and so is its epistemic increment if the model target is $\eta^{\mathrm L}$. It varies at the organism's timescale, so predictability is not what excludes it.
\end{enumerate}
The ligand's chemical potential couples to the blanket above thermal noise, is predictable at the cell's horizon and enters its preferences; gravity fails the coupling clause, with a signal-to-noise ratio more than thirteen orders of magnitude below one per particle and reading, against a chemical estimate that integrates about $7.9\times10^3$ molecule counts, and it fails the relevance clause by construction. The two readings do not pool equally: the count $7.9\times10^3$ is $5\pi D_{\mathrm m}r_{\mathrm{cell}}\bar c_{\mathrm L}T_{\mathrm{int}}/3$, the inverse relative variance in the leading term, and with the exact variance the chemical reading averages about 315 ligand molecules in the ball over about 31 effective snapshots, some $10^4$ molecule-samples; a gravity reading that pooled as many independent protein heights would have ratio $2\times10^{-10}$, still nearly ten orders of magnitude below one. Item (4) is a known physical bound recovered: its number is the barometric exponent $x_{\mathrm g}$ of equilibrium statistical mechanics entered into a Fisher-information identity, no element of the active-inference framework changes it, and the framework contributes only its classification under the coupling clause of \cref{prop:env-selection}. \Cref{fig:env-cell} draws items (1) and (4) against the radius of the buoyant particle: below $x_{\mathrm g}=1$ the ratio of gravity sensing stays under $1/12$, below both ligand signals, and the exact ratio $1-(x_{\mathrm g}/2)^2/\sinh^2(x_{\mathrm g}/2)$ passes the ratio $0.79$ of the 1 percent signal at $x_{\mathrm g}=4.66$ and approaches one from below where sedimentation dominates, so that the small-$x_{\mathrm g}$ form $x_{\mathrm g}^2/12$ overstates it there.
\end{example}
% src: H5-E1 with H5-R4; buoyant mass per critic 2.6

\begin{remark}[Interpretation: neurons, statocytes and vestibular cells]\label{rem:env-neurons}
This remark is an interpretation, and its physiological facts are textbook statements that we did not check against primary sources. A neuron in a body does not sense gravity because a molecular sensor has the negligible coupling computed for the protein in item (4) of \cref{ex:env-cell}, and because its preferred outcomes, transmitter release and firing, do not depend on the direction of gravity. Cells that do sense gravity use dense bodies of the size and density excess of the sedimenting particle of that item, the statoliths of plants and the otoconia of the vestibular organs, for which one height reading has a ratio of one and gives the sign of the axial acceleration with negligible error, and serve organisms whose preferred outcomes, a growth direction or balance, depend on that direction. Whether a variable belongs to an agent's environment therefore depends on the agent's scale and preferences, and an organism, a meta-agent of cells, senses gravity through specialized members. The statement that nerve cells sense chemical potentials but not gravity holds for typical neurons in this sense and fails for vestibular hair cells, whose transduction is mechanical.
\end{remark}
% src: H5-R5 (INTERPRETATION)

\subsection{Aggregates as external states}\label{sup:env-aggregates-as-external-states}

\begin{definition}[Aggregates as external states]\label{def:env-aggregates}
Let a population have joint state $x_t$, the agents' states and the world, with linear dynamics $x_{t+1}=\mathbb Ax_t+\mathfrak n_t$ and $\mathfrak n_t$ independent across rounds. A declared aggregate is $\mathsf Lx_t$, with $\mathsf L$ of full row rank fixed in advance by an institution or a measurement procedure. It is exactly closed when $\ker\mathsf L$ is invariant under $\mathbb A$, so that $\mathsf Lx_{t+1}=\bar{\mathbb A}\mathsf Lx_t+\mathsf L\mathfrak n_t$ for some $\bar{\mathbb A}$ and every initial law (\cref{thm:met-abstract-physics} without input); closed of order $d$ when $d$ further linear functionals $\mathsf L'x_t$ can be adjoined so that $(\mathsf Lx_t,\mathsf L'x_t)$ is exactly closed, the least such $d$ being its minimal order $d^\ast$ in the sense of \cref{def:met-abstract-physics}(iii), which counts adjoined functionals of the joint state and bounds from above the tolerance-dependent closure order $\mathrm{ord}(\alpha_{\mathrm M})$ of \cref{def:agy-closure}(iii), a count of functionals of the filter state (see the discussion after \cref{thm:met-abstract-physics}); and approximately closed when its memory is below a declared bound. For a scalar stationary aggregate $\bar x_t$ the lag-$u$ Markov error is
\[
\Delta^{\mathrm{agg}}_u=\bigl\lvert\Corr(\bar x_t,\bar x_{t+u})-\Corr(\bar x_t,\bar x_{t+1})^u\bigr\rvert,
\]
which compares with the lag-one correlation, unlike the Markov error $\Xi$ of \cref{thm:agy-closure-average}, which compares with the slowest eigenvalue.
\end{definition}
% src: H5-D3 (closure vocabulary of decision 19; Xi^agg renamed Delta^agg_u, critic 2.2, 2.3)

\begin{proposition}[Exact closure of aggregates]\label{prop:env-aggregate-closure}
\begin{enumerate}
\item[(a)] (Polls and opinion averages.) For a linear recursion $x_{t+1}=Wx_t+\zeta_t$ on the agents, receiver first, the average $\bar x=N^{-1}\one^\top x$ is exactly closed if and only if all column sums of $W$ are equal, $\one^\top W=\bar w\one^\top$; then $\bar x_{t+1}=\bar w\bar x_t+N^{-1}\one^\top\zeta_t$. This is \cref{def:met-lumpable} for the partition with one block: every agent carries the same total weight in all agents' updates, its own included.
\item[(b)] (Interchangeable types.) Let the agents form types $\mathfrak t=1,\dots,\mathfrak T$ of sizes $N_{\mathfrak t}$, with
\[
x_{i,t+1}=A_{\mathfrak t}x_{i,t}+\sum_{\mathfrak t'}\mathsf B_{\mathfrak t\mathfrak t'}\bar x_{\mathfrak t',t}+\mathsf H_{\mathfrak t}\eta_t+w^{\mathrm{com}}_{\mathfrak t,t}+w^{\mathrm{id}}_{i,t}\qquad(i\in\mathfrak t),
\]
where $\bar x_{\mathfrak t}$ are the type averages, $w^{\mathrm{com}}$ common type shocks and $w^{\mathrm{id}}_{i,t}\sim\Ncal(0,Q_{\mathrm{id}})$ idiosyncratic shocks, all independent, and $\eta$ is a Gauss--Markov world. Then $(\eta_t,\bar x_{1,t},\dots,\bar x_{\mathfrak T,t})$ is exactly closed, with $\bar x_{\mathfrak t,t+1}=A_{\mathfrak t}\bar x_{\mathfrak t,t}+\sum_{\mathfrak t'}\mathsf B_{\mathfrak t\mathfrak t'}\bar x_{\mathfrak t',t}+\mathsf H_{\mathfrak t}\eta_t+w^{\mathrm{com}}_{\mathfrak t,t}+\bar w^{\mathrm{id}}_{\mathfrak t,t}$ and $\bar w^{\mathrm{id}}_{\mathfrak t,t}\sim\Ncal(0,Q_{\mathrm{id}}/N_{\mathfrak t})$. If the type-level transition $\mathbb A_{\mathrm{ty}}=\diag(A_{\mathfrak t})+[\mathsf B_{\mathfrak t\mathfrak t'}]$ has $\rho(\mathbb A_{\mathrm{ty}})<1$, the stationary deviation of the type averages from their infinite-population limit, the same dynamics without $\bar w^{\mathrm{id}}$, is independent of that limit, with covariance
\[
\Sigma^{\mathrm{fl}}=\sum_{u\ge0}\mathbb A_{\mathrm{ty}}^u\diag\bigl(Q_{\mathrm{id}}/N_{\mathfrak t}\bigr)(\mathbb A_{\mathrm{ty}}^u)^\top,\qquad\Vert\Sigma^{\mathrm{fl}}\Vert\le\frac{\Vert Q_{\mathrm{id}}\Vert}{\min_{\mathfrak t}N_{\mathfrak t}}\sum_{u\ge0}\Vert\mathbb A_{\mathrm{ty}}^u\Vert^2.
\]
\item[(c)] If in (b) two agents $i,i'$ of one type have own dynamics $A_i\ne A_{i'}$ in place of the common $A_{\mathfrak t}$, the type averages are not exactly closed.
\end{enumerate}
\end{proposition}
% src: H5-P5 (one formula for the block-sum condition, def:met-lumpable, critic V6)

Proof in \cref{app:env}; checked in \cref{chk:env-aggregates}. For scalar states, (b) is \cref{def:met-lumpable} for the partition into types with the world as one more block, and the law of large numbers adds an error whose variance is of order $1/\min_{\mathfrak t}N_{\mathfrak t}$, so that the deviation itself is of order $(\min_{\mathfrak t}N_{\mathfrak t})^{-1/2}$. Exact closure of an aggregate is the linear analog of lumpability of a partition of a Markov chain \citep{kemeny1960finite}, which we know in restatement, and type homogeneity is its population form.

\begin{proposition}[Heterogeneous timescales give aggregates memory]\label{prop:env-aggregate-memory}
\begin{enumerate}
\item[(a)] A stationary Gaussian scalar process with positive variance is Markov if and only if its autocovariance is $\mathsf R(u)=\mathsf R(0)\alpha_\ast^{\lvert u\rvert}$ for some $\alpha_\ast$; a scalar aggregate is therefore Markov exactly when $\Delta^{\mathrm{agg}}_u=0$ for every $u$.
\end{enumerate}
Let $x_{i,t+1}=\alpha_ix_{i,t}+w^{\mathrm{com}}_t+w^{\mathrm{id}}_{i,t}$ for $i\le N$, with persistences $\alpha_i\in(-1,1)$, a common shock of variance $q_{\mathrm{com}}$ and idiosyncratic shocks of variance $q_{\mathrm{id}}$, all independent, $q_{\mathrm{com}}+q_{\mathrm{id}}>0$; let $\mathbb Q$ be the empirical law of the $\alpha_i$ and $\bar x_t$ the stationary average.
\begin{enumerate}
\item[(b)] $\mathsf R(u)=\Cov(\bar x_t,\bar x_{t+u})=\int\alpha^u\mathbb M(\mathrm d\alpha)$ for $u\ge0$, with the positive measure
\[
\mathbb M(\mathrm d\alpha')=q_{\mathrm{com}}\Bigl(\int\frac{\mathbb Q(\mathrm d\alpha)}{1-\alpha\alpha'}\Bigr)\mathbb Q(\mathrm d\alpha')+\frac{q_{\mathrm{id}}}{N}\frac{\mathbb Q(\mathrm d\alpha')}{1-\alpha'^2}.
\]
\item[(c)] $\bar x$ is Markov if and only if all $\alpha_i$ are equal.
\item[(d)] If the $\alpha_i$ lie in $[\alpha_-,\alpha_+]\subset[0,1)$, then, with $\bar{\mathbb M}=\mathbb M/\mathbb M(\R)$, for every $u$
\[
\Delta^{\mathrm{agg}}_u\le\frac{\Var_{\bar{\mathbb M}}(\alpha)}{(1-\alpha_+)^3}\le\frac{(\alpha_+-\alpha_-)^2}{4(1-\alpha_+)^3}.
\]
\item[(e)] If the $\alpha_i$ take $\mathfrak T$ distinct values, the type averages are exactly closed and the minimal order of $\bar x$ is $d^\ast=\mathfrak T-1$.
\item[(f)] (Long memory.) As $N\to\infty$ with the common shock only, let $\mathbb Q$ have density $(1-\alpha)^{\mathfrak p-1}\mathsf f_0(\alpha)$ on $[0,1)$, with $\mathsf f_0$ continuous on $[0,1]$, $\mathsf f_0(1)>0$ and $\frac12<\mathfrak p<1$. Then $\mathfrak m_l=\int\alpha^l\mathbb Q(\mathrm d\alpha)\sim\mathsf f_0(1)\Gamma(\mathfrak p)l^{-\mathfrak p}$, and the limiting aggregate $\bar x_t=\sum_{l\ge0}\mathfrak m_lw^{\mathrm{com}}_{t-1-l}$ has
\[
\mathsf R(u)\sim q_{\mathrm{com}}\bigl(\mathsf f_0(1)\Gamma(\mathfrak p)\bigr)^2\mathrm B(1-\mathfrak p,2\mathfrak p-1)u^{1-2\mathfrak p},
\]
with $\Gamma$ the Gamma function and $\mathrm B$ the Beta function, so $\sum_u\mathsf R(u)=\infty$: the aggregate has long memory although every agent is an autoregression of order one.
\end{enumerate}
\end{proposition}
% src: H5-P6 and the lemma of H5-D3; critic V6

Proof in \cref{app:env}; checked in \cref{chk:env-aggregates}. Parts (c) to (f) are \cref{asm:int-a7} in aggregate form: an aggregate of heterogeneous members is not Markov, the memory it needs grows with the number of distinct timescales it mixes, and with a continuum of timescales reaching the unit root no finite memory suffices. Part (f) is a common-innovation case of the aggregation route to long memory and is not new. The review of \citet{leipus2013aggregation}, whose first two sections we read, traces that route to \citet{granger1980long}, a paper we know from its record and from the review, for autoregressions with a random persistence that has a Beta law. For a mixing density proportional to $(1-\alpha)^{\beta}$ near one and common innovations of infinite variance, the review reports that the average normalized by $1/N$ converges to a moving average in the common innovations with coefficients $\E\alpha^l\sim\mathrm{const}\cdot l^{-\beta-1}$, which are not summable for $\beta<0$; with $\beta=\mathfrak p-1$ these are the weights $\mathfrak m_l$ of (f). The review credits this limit to \citet{puplinskaite2009aggregation}, which we did not open, and notes that the general aggregation scheme of \citet{zaffaroni2004contemporaneous}, which we did not open either, includes the case of common innovations. Part (f) is the Gaussian counterpart with finite variance, for which we give a self-contained proof with the explicit constant. \Cref{ex:met-traders} shows the input-driven form of the same effect: the memory kernel of an average of traders with heterogeneous persistences survives the law of large numbers. The third closure mechanism, timescale separation, is \cref{thm:agy-closure-average}: an outside agent that reads the average of a collective inherits its bound on the Markov error $\Xi$, set by a dispersion of the members' leaks or of the row and column sums of the drift against the gap between the inward and outward rates of the precision-based contrast (reversible drift) or of the drift-based contrast (general drift). Part (d) is the discrete-time analog for uncoupled members, with the dispersion of the persistences against the gap $1-\alpha_+$.

\subsection{Firms and people}\label{sup:env-firms-and-people}

\begin{example}[A firm with a production lag reads one market aggregate]\label{ex:env-firm}
The market state $\eta_t=(\eta_{1,t},\eta_{2,t})$ has independent scalar autoregressive components with persistences $F_1>F_2>0$, a slow trend in demand and a fast shock, and stationary variances $P_{0,1}$ and $P_{0,2}$; the value of the good is $\eta^{\mathrm{val}}_t=\eta_{1,t}+\eta_{2,t}$. At the fully revealing price of \cref{prop:abs-walras}(b) we model the spot price as $p_t=\eta^{\mathrm{val}}_t+\mathsf n^{\mathrm p}_t$, with $\mathsf n^{\mathrm p}_t$ independent of the market state and of variance $(\sum_i1/R^s_i)^{-1}$, the noise of the statistic that the price reveals, of order $1/N$ for $N$ traders. A firm sells at $t+h$ what it produces at $t$, with profit $p_{t+h}a-\frac{\mathsf k_a}2a^2$ for output $a$; its optimal output is $\E[p_{t+h}\mid\text{readings}]/\mathsf k_a$, and its excess loss over full information on $\eta_t$ is $\Var(\E[p_{t+h}\mid\eta_t]\mid\text{readings})/(2\mathsf k_a)$. The decision functional of the present, $(F^h)^\top\one=(F_1^h,F_2^h)$, turns from $(1,1)$ toward $(1,0)$ as $h$ grows. The firm reads one aggregate, which we take to be noiseless: the spot price, functional $(1,1)$, neglecting the price noise $\mathsf n^{\mathrm p}_t$, whose variance is of order $1/N$ (with it, $\mathcal L_{\mathrm{spot}}$ is the formula in \cref{app:env}, and the comparison with the slow index below is that of the noiseless limit); a slow index such as a trend estimate or a long-dated contract, functional $(1,0)$; or a futures price of maturity $h$, the market forecast of the value at $t+h$, functional $(F_1^h,F_2^h)$. Up to the factor $1/(2\mathsf k_a)$ the excess losses are
\[
\mathcal L_{\mathrm{spot}}(h)=(F_1^h-F_2^h)^2\frac{P_{0,1}P_{0,2}}{P_{0,1}+P_{0,2}},\qquad\mathcal L_{\mathrm{slow}}(h)=F_2^{2h}P_{0,2},\qquad\mathcal L_{\mathrm{fut}}(h)=0,
\]
so the futures price at the firm's own lag is sufficient. This holds by definition of the futures functional, for any dynamics, since the futures price is the conditional expectation that the output decision needs and a conditional expectation is sufficient for itself; \cref{thm:env-hayek-horizon}(a) with $M=(F_1^h,F_2^h)$ states the same fact. That the futures price of the firm's own lag suffices is the separation theorem of production under price uncertainty, here in its risk-neutral form and by definition of the futures functional: for a competitive firm that can trade futures, production decisions depend only on the futures price and input costs, whatever the firm's subjective distribution of the future spot price \citep{feder1980futures}, which we know from its record and abstract. The sufficiency of the futures price is therefore a known result recovered, and only the comparison of the other two readings is specific to the declared two-mode exterior. The spot price beats the slow index exactly for $h<h^\times$, with
\[
h^\times=\frac{\ln\bigl(1+(1+P_{0,2}/P_{0,1})^{1/2}\bigr)}{\ln(F_1/F_2)}.
\]
For $F=\diag(0.95,0.6)$ and $P_0=\diag(1,2)$, $h^\times=2.187$, and at $h=4$ the information-optimal rank-one reading of the market state is the slow component $(1,0)$, while the decision-optimal reading is the futures functional $(1,0.159)$ (\cref{chk:env-firm}): the canonical variate is what to read for prediction in general, and the decision functional what to read for this decision. The price is treated as exogenous, without price impact.
\end{example}
% src: H5-E2; critic V12 (worked example, price exogenous); derivation in app:env

\begin{remark}[Interpretation: the firm as a meta-agent]\label{rem:env-firm-meta-agent}
This remark is an interpretation. Read as a meta-agent (\cref{def:met-coarse}), the firm has as sensory states its members' readings of aggregates (prices, indices, futures, its own sales), as active states its output and prices, which enter the market as performative acts (\cref{prop:fw-performative}), as internal states a belief about the market pooled by consensus on information contributions (\cref{prop:agy-consensus}), and as external states the market state and the other firms' aggregate output. Every exterior variable it reads is an aggregate, closed or carrying memory as \cref{prop:env-aggregate-closure,prop:env-aggregate-memory} say, and its horizon and decision functionals fix which aggregates it should read. Whether its pooled belief and output meet the clauses of \cref{thm:met-h5a} is the subject of \cref{sec:met}, and \cref{ex:met-two-firms} computes them for two firms whose abstract physics contains the competitor's coarse states. A firm whose output is a non-negligible share of supply and that reads the price as a message about demand without removing its own act misreads its own act as evidence (\cref{prop:fw-performative}(a)); removing it is the cavity rule of \cref{thm:hie-messages}, and the size of the error depends on the price impact, which we have not modeled.
\end{remark}
% src: H5-R7 (INTERPRETATION; interfaces IF-a, IF-b)

\begin{example}[A person senses the economy through one or two numbers]\label{ex:env-person}
The economy's state has four independent modes with correlation times $\tau=(40,8,2,0.5)$ periods and innovation variances $(0.05,0.2,0.5,1)$, so stationary variances $(1.03,0.90,0.79,1.02)$, read in mode coordinates as trend, cycle, local prices and noise. A person has three equally weighted decisions $\mathfrak d$ with quadratic losses: consumption at horizon 1 with functional $(0.3,0.3,1,0.6)$, a job decision at horizon 8 with $(1,1,0.3,0)$, and saving at horizon 40 with $(1,0.4,0,0)$. Its total excess loss is $\tr\mathsf J-\tr(\mathsf G_s\mathsf J)$ with $\mathsf J=\sum_{\mathfrak d}P_0^{1/2}(F^{h_{\mathfrak d}})^\top\psi_{\mathfrak d}\psi_{\mathfrak d}^\top F^{h_{\mathfrak d}}P_0^{1/2}$, so under the Shannon budget $I(s;\eta_t)\le I_{\max}$ of rational inattention the optimal signal reads the leading eigenvectors of $\mathsf J$ by reverse water-filling (\cref{thm:env-canonical}(c)), at most one channel per decision. The eigenvalues of $\mathsf J$ are $(1.118,0.263,0.016,0)$. The first channel reads the functional $(1,0.46,0.31,0.04)$ on (trend, cycle, local prices, noise), a slow aggregate, and the second $(0.31,-0.15,-1,-0.13)$, mainly local prices. With $0.05$ or $0.3$ nats per period the person reads the slow aggregate alone; the second channel enters above $0.72$ nats, and the third, mainly the cycle, only above $3.50$ nats (\cref{chk:env-person}). The economy, for this person, is one or two numbers, chosen by the horizons and weights of its own decisions.
\end{example}
% src: H5-E3; activation budgets recomputed by the editor

\subsection{Sensing across levels and memberships}\label{sec:env-levels}

A person reads aggregates that meta-agents produce from their members' sensing, and belongs to several meta-agents at once. Two results bound what passes through such chains and how one budget is shared among memberships.

\begin{lemma}[Composition across levels]\label{lem:env-composition}
If $\mathsf X_1\to\mathsf X_2\to\mathsf X_3$ is a jointly Gaussian Markov chain with nonsingular covariances and $\mathsf C_{12}$, $\mathsf C_{23}$ and $\mathsf C_{13}$ are the whitened cross-covariances of $(\mathsf X_1,\mathsf X_2)$, $(\mathsf X_2,\mathsf X_3)$ and $(\mathsf X_1,\mathsf X_3)$, then $\mathsf C_{13}=\mathsf C_{12}\mathsf C_{23}$. Hence $\sigma_{\mathsf j}(\mathsf C_{13})\le\sigma_{\mathsf j}(\mathsf C_{12})\sigma_1(\mathsf C_{23})$ and $\rank\mathsf C_{13}\le\min(\rank\mathsf C_{12},\rank\mathsf C_{23})$.
\end{lemma}
% src: H5-L5

Proof in \cref{app:env}; checked in \cref{chk:env-composition}. Along a chain from the world through members' sensing and a meta-agent's aggregate to a person's reading of it, every canonical correlation between the world and the reading is at most the product of the stages' leading canonical correlations, and a broadcast code $T_c$ of rank $r_c$ passes at most $r_c$ canonical variates; this is consistent with \cref{conj:sca-rank-down} and does not prove it, since that conjecture concerns hop distance. If a meta-agent's sensory statistic is linear in its members' readings (\cref{def:met-coarse}), its exterior channels and their noises fix its budget, and its closure timescale sets its horizon, of order $1/\varpi_1$ for the average of \cref{thm:agy-closure-average}, then \cref{thm:env-canonical,prop:env-horizon,prop:env-diffusive} say which coarse exterior variables it should sense: the leading canonical variates at its horizon, slower and coarser for slower meta-agents. Whether a collective has such a statistic and horizon is the subject of \cref{sec:met}.

\begin{proposition}[A person with physical and abstract channels]\label{prop:env-multiplex}
Let a person's decisions be $\mathsf Y=(\psi_{\mathfrak d}^\top\eta_{t+h_{\mathfrak d}})_{\mathfrak d}$, let it read the physical world through $s^{\mathrm{ph}}=H^{\mathrm{ph}}\eta_t+v^{\mathrm{ph}}$ and the aggregates of its memberships, one modality $c\in\mathcal M_i\subseteq\mathcal M$ each, through $y^{\mathrm{agg}}=(T^{\mathrm{agg}}_c\eta_t+\nu_c)_c$, all noises independent, and put $\mathsf Z_h=\bigl[(F^{h_{\mathfrak d}})^\top\psi_{\mathfrak d}\bigr]_{\mathfrak d}$.
\begin{enumerate}
\item[(a)] $I(\mathsf Y;s^{\mathrm{ph}},y^{\mathrm{agg}})=I(\mathsf Y;y^{\mathrm{agg}})+I(\mathsf Y;s^{\mathrm{ph}}\mid y^{\mathrm{agg}})$, each term in closed form through conditional covariances.
\item[(b)] (Sufficiency.) $I(\mathsf Y;s^{\mathrm{ph}}\mid y^{\mathrm{agg}})=0$ if and only if $\mathsf Z_h^\top\Cov(\eta_t\mid y^{\mathrm{agg}})(H^{\mathrm{ph}})^\top=0$. With noiseless aggregates, $\operatorname{col}\mathsf Z_h\subseteq\operatorname{row}T^{\mathrm{agg}}$ suffices, which is the condition of \cref{thm:env-hayek-horizon}(a).
\item[(c)] (Long horizons.) For a scalar decision $\mathsf Y=\psi^\top\eta_{t+h}$ and every physical channel, $I(\mathsf Y;s^{\mathrm{ph}}\mid y^{\mathrm{agg}})\le-\frac12\ln(1-\mathfrak r_h^2)$ with $\mathfrak r_h^2=\Var(\psi^\top F^h\eta_t\mid y^{\mathrm{agg}})/\Var(\mathsf Y\mid y^{\mathrm{agg}})$. In a reversible exterior with noiseless aggregates that read the $r_{\mathrm{agg}}$ slowest mode amplitudes, for $h\ge1$ and $\psi\ne0$,
\[
\mathfrak r_h^2\le\frac{\lvert\lambda_{r_{\mathrm{agg}}+1}(F)\rvert^{2h}\sum_{\mathsf j>r_{\mathrm{agg}}}(\mathsf e_{\mathsf j}^\top P_0^{1/2}\psi)^2}{\psi^\top\bigl(P_0-F^hP_0(F^h)^\top\bigr)\psi},
\]
which decays geometrically in $h$.
\item[(d)] (Attention across memberships.) If $F$ and $Q$ are block diagonal across contents, as in \cref{prop:abs-decoupling}(a), and each decision loads on one content, the decision matrix $\mathsf J$ of \cref{ex:env-person} is block diagonal across memberships. The budget-optimal sensor then reads channels within one membership's content each, when eigenvalues of different blocks differ, and the allocation across memberships is coupled only through the common water level: raising the decision eigenvalues of one membership leaves the capacity spent on every other membership's channels unchanged or lowers it.
\end{enumerate}
\end{proposition}
% src: H5-P7

Proof in \cref{app:env}; checked in \cref{chk:env-multiplex}. For long-horizon decisions aggregates of the slow modes are therefore sufficient up to an explicit error, and a physical channel can add much only for short-horizon decisions. Part (d) answers the single-agent part of the open problem of attention across modalities (\cref{open:abs-attention}) for a Shannon budget and pragmatic value alone; the case of a channel count $\kappa$ and of the full objective $G^{\mathrm{att}}_h$ of \cref{def:env-attention-objective} remains open. In the case that (d) covers, with block-diagonal decision structure one budget shared across modalities does not mix them within channels, but it couples the amounts of attention through the water level. With decision eigenvalues $(2.0,0.4)$ for a firm and $(1.2,0.9)$ for a polity, a person attends only the firm's leading channel below $0.26$ nats, adds the polity's two channels at $0.26$ and $0.54$ nats, and the firm's second channel only above $1.76$ nats.

\begin{remark}[Reading an aggregate one contributes to]\label{rem:env-own-aggregate}
When a person's own report enters an aggregate it reads, a poll it answered or a price at which it trades, a reader at interpretation level 0 counts its own information twice: in market form, by \cref{prop:abs-walras}(c), its believed precision exceeds the true one by the precision of its prior and of its own signal. \Cref{lem:env-composition} assumes that the reading is a downstream stage of the chain; with own contributions the chain is a loop, and the correct rule removes the own contribution before reading, the cavity rule of \cref{thm:hie-messages}. Literal reading of such an aggregate is the overvaluation case of \cref{sec:mis}, treated as literal top-down reading in \cref{prop:hie-literal-once}.
\end{remark}
% src: H5-R8 (interface IF-b)

\begin{remark}[Interpretation: corporations and people sensing abstract environments]\label{rem:env-interpretation}
This remark is an interpretation of the results of \cref{sec:env} and of this section, which make the author's examples precise in linear-Gaussian models for sensors optimal under the declared objective of \cref{def:env-attention-objective}; the one-step expected free energy of the framework does not imply that objective (\cref{rem:env-efe-horizon-free,rem:env-two-step}), so each sentence below says what such a sensor reads and not what every agent of the framework senses. A corporation senses an informational environment of prices, indices and contracts because at its decision horizon the variables that are predictive and relevant to its acts are slow aggregates, which under the homogeneity conditions of \cref{prop:env-aggregate-closure} are closed processes that it can model as its exterior (\cref{thm:env-hayek-horizon,ex:env-firm}). A person senses the economy through one or two numbers when its long-horizon decisions load on slow modes and its attention is limited (\cref{ex:env-person}), and senses the physical world for its short-horizon decisions (\cref{prop:env-multiplex}(c)). A cell senses its ligand's chemical potential and not gravity because of coupling, horizon and relevance together (\cref{ex:env-cell}). In each case an abstract environment is the exterior seen through the leading canonical variates and decision functionals at the agent's horizon. \Citet{fields2022competency} propose a related scale-free view, competent behavior in metabolic, transcriptional, morphological and spatial spaces framed by hierarchical active inference; \cref{thm:env-canonical,thm:env-control,thm:env-hayek-horizon} give one quantitative criterion for which space an agent behaves in.
\end{remark}
% src: H5-R9 (INTERPRETATION)

\subsection{Open problems}\label{sec:env-open}

\begin{openproblem}[Sensor design for filtering agents]\label{open:env-filtering}
For an agent that integrates its readings with a Kalman filter, which sensor of $\kappa$ channels maximizes $I(s_{\le t};\eta_{t+h})$ in steady state? \Cref{lem:env-markov} settles summaries of a fully observed past, and \cref{thm:env-canonical} applied to a window of past observations settles a fixed window. With memory a sensor of one variable also informs others through the dynamics, so the answer is not the canonical variates in general: a single channel that reads a mixture of a slow and a fast mode yields a process whose past carries information about both. The problem couples the choice of sensor with the Riccati equation.
\end{openproblem}
% src: H5-O1

\begin{conjecture}[Selection beyond linear-Gaussian exteriors]\label{conj:env-nonlinear}
For a stationary exterior with finite predictive information and an agent with a budget on $I(s;\eta_t)$, the optimal representations for prediction at horizon $h$, the information bottleneck between present and future, concentrate as $h$ grows on functions of the exterior that are invariant under the slowest part of the dynamics: eigenfunctions of the transfer operator of largest modulus, or their nonlinear analogs.
\end{conjecture}
% src: H5-C1

For linear-Gaussian exteriors with diagonalizable $F$ and a gap in the moduli of its eigenvalues, the statement of \cref{conj:env-nonlinear} follows from \cref{thm:env-canonical}(b) and \cref{prop:env-horizon}(b): the budgeted optimum reads the leading canonical variates of $(\eta_t,\eta_{t+h})$, and along the horizons at which it reads $\mathsf m$ of them, if $\lvert\lambda_{\mathsf m}(F)\rvert>\lvert\lambda_{\mathsf m+1}(F)\rvert$, the bound of \cref{prop:env-horizon}(b) with $\kappa=\mathsf m$ places their span within an angle of order $(\lvert\lambda_{\mathsf m+1}(F)\rvert/\lvert\lambda_{\mathsf m}(F)\rvert)^h$ of the span of the $\mathsf m$ slowest mode amplitudes. Defective $F$ and counts $\mathsf m$ without such a gap are not covered; discrete and nonlinear versions, such as the renormalizing generative models of \citet{friston2025pixels}, are covered by no proof here. For field theories, the degrees of freedom that information-bottleneck compression keeps are reported to be the operators of lowest scaling dimension \citep{gordon2021relevance}, a result we know from its abstract that would tie the conjecture to the relevant variables of \cref{sec:ren}. Three narrower questions remain open: whether the number of variables read at a fixed budget is non-increasing in $h$ for non-reversible exteriors, where transient growth of a non-normal $F^h$ can raise a canonical correlation at short horizons, so that we expect counterexamples, although we did not search for them; whether, in a network whose agents' decision eigenvalues depend on the state of each modality, the water-level coupling of \cref{prop:env-multiplex}(d) generates a cross-modal coupling $a_\times$ in the sense of \cref{prop:abs-crossover}, the network part of the open problem of attention across modalities in \cref{sec:abs}, which stays open together with its single-agent case for a channel count $\kappa$ and the full objective $G^{\mathrm{att}}_h$; and which sensor model fits a given cell, since absorbing spheres and receptors with binding kinetics have constants other than those of \cref{prop:env-concentration} \citep{bialek2005physical}.
